\chapter{Cohomología del acarreo helicoidal: evaluación de cociclos y dimensión mínima de la memoria acumulativa}
\label{chap:elevacion-cohomologica-dimension-rev2}
\index{dimensión!rango cohomológico}
\index{acarreo!integral discreta}
\index{memoria!elevación dimensional}

Una coordenada periódica registra la posición visible; el acarreo registra
cuántas veces la trayectoria ha atravesado su frontera. La primera retorna,
la segunda se acumula. Esta diferencia convierte la periodicidad en una
elevación: dos instantes pueden compartir fase visible y pertenecer a estados
distintos porque conservan números de vuelta diferentes.

En el alfabeto nonádico, para \(n\in\mathbb Z_{\geq1}\), definamos
\[
 r(n)=1+((n-1)\bmod9),
 \qquad
 q(n)=\left\lfloor\frac{n-1}{9}\right\rfloor,
\]
y el levantamiento helicoidal
\begin{equation}
 H_9(n)=(r(n),q(n)).
 \label{eq:elevacion-helicoidal-nonadica-rev2}
\end{equation}
Entonces
\[
 H_9(n+9)=H_9(n)+(0,1):
\]
la fase retorna y la memoria cambia. Si \(c_k\in\mathbb Z\) es el acarreo
del paso \(k\to k+1\), su integral discreta
\begin{equation}
 Q_N=Q_0+\sum_{k=0}^{N-1}c_k
 \label{eq:integral-discreta-carry-rev2}
\end{equation}
constituye la coordenada acumulativa del levantamiento. Su independencia es
algebraica; una realización métrica posterior puede representarla mediante
una dirección transversal u ortogonal.

\section{Cociclos independientes y número mínimo de coordenadas de memoria}

Sea \(G\) el grafo orientado de estados visibles y sean
\(c_1,\ldots,c_d\in Z^1(G;\mathbb Z)\) cociclos enteros. Para una ruta
\(\gamma\), definamos
\[
 Q_j(\gamma)=\sum_{e\in\gamma}c_j(e),
 \qquad
 \widetilde\gamma=
 \bigl(q_{\rm vis}(\gamma),Q_1(\gamma),\ldots,Q_d(\gamma)\bigr).
\]
Sobre una ruta abierta esta evaluación retiene el representante del cociclo
y los extremos. En efecto, con la convención
\((\mathrm d\phi)(e)=\phi(t(e))-\phi(s(e))\),
\begin{equation}
 Q_j^{\,c_j+\mathrm d\phi}(\gamma)
 =Q_j^{\,c_j}(\gamma)+\phi(t\gamma)-\phi(s\gamma).
 \label{eq:cambio-representante-cociclo-ruta-rev2}
\end{equation}
Por tanto, para una ruta abierta se fijan el representante y sus extremos,
y la memoria conserva esos datos. Si \(\gamma\) es un ciclo, el término de
frontera se anula y \(Q_j(\gamma)\) depende únicamente de la clase
\([c_j]\in H^1(G;\mathbb Z)\).

\begin{theorem}[Criterio cohomológico de rango dimensional]
\label{thm:rango-dimensional-cohomologico-rev2}
Supóngase que las clases de \(c_1,\ldots,c_d\) generan un submódulo libre de
rango \(d\) en \(H^1(G;\mathbb Z)\). Sea
\[
 Q:Z_1(G;\mathbb Z)\longrightarrow\mathbb Z^d
\]
su evaluación sobre ciclos. Si una familia de \(m\) contadores aditivos se
obtiene mediante un homomorfismo \(A:\mathbb Z^d\to\mathbb Z^m\), de modo
que
\[
 \widehat Q:=A\circ Q:Z_1(G;\mathbb Z)\longrightarrow\mathbb Z^m,
\]
y permite reconstruir las evaluaciones, es decir, existe
\(B:\mathbb Z^m\to\mathbb Z^d\) con
\(Q=B\circ\widehat Q\), entonces
\[
 \boxed{m\geq d}.
\]
Por tanto, el levantamiento necesita al menos \(d\) coordenadas de memoria
independientes además de la coordenada visible.
\end{theorem}

\begin{proof}
El emparejamiento
\(H^1(G;\mathbb Z)\times H_1(G;\mathbb Z)\to\mathbb Z\) y la independencia
de \([c_1],\ldots,[c_d]\) implican
\(\operatorname{rk}\operatorname{im}Q=d\). Como
\(Q=B\circ\widehat Q\),
\[
 d=\operatorname{rk}\operatorname{im}Q
 \leq\operatorname{rk}\operatorname{im}\widehat Q\leq m.
\]
Ningún homomorfismo hacia menos de \(d\) contadores enteros admite una
inversa aditiva sobre el submódulo generado por las evaluaciones \(Q_j\).
La condición más fuerte \(B\circ A=I_{\mathbb Z^d}\) es suficiente, pero la
cota sólo necesita la reconstrucción de \(Q\) sobre las evaluaciones
realizadas por los ciclos.
\end{proof}

La dimensión genealógica cuantifica aquí grados de libertad necesarios para
recuperar el transporte. El teorema no identifica todavía ese rango con una
dimensión espacial física: proporciona el dominio mínimo sobre el que una
realización métrica puede actuar sin borrar los cociclos independientes.

\section{Torre nonádica y conservación proyectiva de la unidad}

Sea
\[
 \mathcal C_9^{(d)}=(\mathbb Z/9\mathbb Z)^d.
\]
Para una marca \(\star\in\mathbb Z/9\mathbb Z\), definimos la inclusión
\[
 h_{d,\star}:\mathcal C_9^{(d)}\longrightarrow
 \mathcal C_9^{(d+1)},
 \qquad h_{d,\star}(x)=(x,\star),
\]
y la proyección
\[
 \pi_{d+1,d}(x_1,\ldots,x_{d+1})=(x_1,\ldots,x_d).
\]

\begin{proposition}[Conservación por elevación dimensional]
\label{prop:conservacion-elevacion-dimensional-rev2}
Se cumple
\[
 \pi_{d+1,d}\circ h_{d,\star}=\operatorname{id}_{\mathcal C_9^{(d)}}.
\]
En consecuencia, la celda de dimensión \(d\) reaparece exactamente como una
sección coordenada de la celda de dimensión \(d+1\): una rebanada afín, y un
subgrupo cuando \(\star=0\). El término \emph{cara de codimensión uno} se
reserva para una realización cubical adicional expresamente declarada. Para las medidas
uniformes normalizadas \(\mu_d\),
\[
 (\pi_{d+1,d})_*\mu_{d+1}=\mu_d.
\]
\end{proposition}

\begin{proof}
La primera identidad resulta de eliminar la última coordenada recién
añadida. Para un punto \(x\in\mathcal C_9^{(d)}\), la fibra de
\(\pi_{d+1,d}\) contiene nueve puntos, cada uno con masa
\(9^{-(d+1)}\); su masa total es \(9^{-d}=\mu_d(\{x\})\).
\end{proof}

La primera igualdad expresa conservación estructural; la segunda,
conservación normalizada. Elevar no reemplaza la capa precedente: la incluye
como sección y distribuye la nueva coordenada sobre sus fibras. En dimensión
tres, \(|\mathcal C_9^{(3)}|=729\); su completación decimal y la
descomposición disjunta
\[
 1000=729+243+27+1
\]
permanecen en la residencia de completación cúbica, donde se distinguen
interior, caras, aristas y vértice añadido.

\section{De la coordenada acumulativa a la métrica}

El levantamiento de cociclos y la elección de una métrica son operaciones
sucesivas. El primero construye coordenadas independientes de memoria; la
segunda declara una forma cuadrática que mide sus variaciones. En APP, el
bloque central
\[
 B_c=\begin{pmatrix}7&2\\2&7\end{pmatrix}=9P_++5P_-
\]
produce la normalización de referencia
\[
 M_c:=\frac{B_c}{\sqrt{45}}\in SO^+(1,1),
\]
demostrada en el \cref{prop:app-trit-espectral-markov-lorentz-rev8}.
En la pantalla cúbica, la realización física selecciona la
firma lorentziana correspondiente y el lector temporal de TPK orienta su
dinámica. La cadena causal es
\[
 \begin{aligned}
 \text{retorno con acarreo}
 &\longrightarrow\text{cociclo acumulativo}
 \longrightarrow\text{rango de memoria}\\
 &\longrightarrow\text{forma cuadrática}
 \longrightarrow\text{métrica realizada}.
 \end{aligned}
\]

\paragraph{Separación entre el rango cohomológico interno y su realización espacio--temporal.}
La hélice nonádica y la integral del acarreo proporcionan las premisas del
criterio cohomológico de rango y de su cota mínima. La torre de celdas y la
conservación de la medida se deducen de la proyección finita. La identificación
del rango genealógico con una dimensión espacio-temporal requiere el funtor
métrico declarado en Mecánica Dimensional.
